\documentclass[11pt]{article}

\usepackage[margin=1in]{geometry}
\usepackage{amsmath,amssymb,amsthm}
\usepackage{enumitem}
\usepackage{hyperref}
\usepackage{xcolor}
\usepackage{booktabs}

\hypersetup{colorlinks=true, linkcolor=blue, urlcolor=blue, citecolor=blue}

\title{TOAE209/TOAH209 -- Design and Analysis of Algorithms\\[4pt]
\large Week 5: Theoretical Questions}
\author{Foreign Trade University}
\date{}

\begin{document}
\maketitle

\section*{Instructions}

Part A contains 17 questions requiring written explanations or short proofs.
Part B is a 10-question quiz (true/false and multiple choice). Attempt every
question on your own before consulting Part C (Solutions) at the end of this
document.

\section*{Part A: Theory Questions}

\begin{enumerate}[label=\textbf{A\arabic*.}, leftmargin=2.2em]

\item State the single-source shortest path problem precisely: what is the input
(graph, weights, source), and what is the desired output? Why do we usually
require the edge weights to be non-negative for Dijkstra's algorithm?

\item Describe Dijkstra's algorithm in your own words. What invariant does it
maintain about the set $S$ of ``finalized'' vertices, and what does the relaxation
step do to a tentative distance $d[v]$?

\item Prove that Dijkstra's algorithm is correct on graphs with non-negative
edge weights: when a vertex $u$ is added to $S$, its tentative distance $d[u]$
equals the true shortest-path distance. Identify the exact step in your proof
that uses non-negativity.

\item Analyze the running time of Dijkstra's algorithm implemented with a binary
heap. Count the \textsc{ExtractMin} and \textsc{DecreaseKey}/push operations and
state the overall bound in terms of $n$ and $m$.

\item Give a small directed graph with one negative edge (and no negative cycle)
on which Dijkstra returns an incorrect distance. State the wrong answer Dijkstra
gives and the correct answer, and explain which step of the correctness proof
fails.

\item State the Bellman--Ford algorithm. Why does it perform exactly $n-1$ rounds
of relaxing all edges, and what does the additional ($n$-th) round detect?

\item Compare Dijkstra and Bellman--Ford on (a) the class of graphs they handle
correctly, and (b) their running times. When would you choose each?

\item Define a \emph{spanning tree} and a \emph{minimum spanning tree} of a
connected weighted graph. How many edges does a spanning tree of an $n$-vertex
graph have, and why?

\item Define a \emph{cut} $(S, V \setminus S)$ and what it means for an edge to
\emph{cross} the cut. State the cut property.

\item Prove the cut property: assuming distinct edge weights, the minimum-weight
edge crossing any cut belongs to every MST. Use an exchange argument.

\item State and prove the cycle property: the unique maximum-weight edge on any
cycle belongs to no MST.

\item State Kruskal's algorithm. Prove it correct by showing that every edge it
adds is a minimum-weight crossing edge of some cut (and hence safe by the cut
property).

\item State Prim's algorithm. Prove it correct using the cut property, and
explain in what sense it is the ``graph analogue'' of Dijkstra's algorithm.

\item Describe the union-find (disjoint-set) data structure with \emph{path
compression} and \emph{union by rank}. What are the two operations, and what is
the amortized cost per operation? How is union-find used inside Kruskal's
algorithm?

\item Suppose the edge weights of a connected graph are all \emph{distinct}.
Prove that the minimum spanning tree is unique. What can go wrong (regarding
uniqueness) when weights are allowed to be equal?

\item Using the cut and cycle properties, explain how to decide whether a given
edge $e$ belongs to (i) \emph{every} MST, (ii) \emph{some} but not every MST, or
(iii) \emph{no} MST. Give a characterization for each case.

\item Describe single-link $k$-clustering of maximum spacing and explain its
connection to Kruskal's algorithm. What quantity does ``stopping Kruskal when $k$
components remain'' maximize, and what is the resulting spacing equal to?

\end{enumerate}

\newpage
\section*{Part B: Quiz}

\begin{enumerate}[label=\textbf{B\arabic*.}, leftmargin=2.2em]

\item \textbf{(True/False)} Dijkstra's algorithm always computes correct
shortest-path distances, even when some edge weights are negative (as long as
there is no negative cycle).

\item \textbf{(Multiple choice)} Dijkstra's algorithm implemented with a binary
heap runs in time:
\begin{enumerate}[label=(\alph*)]
    \item $O(n^2)$
    \item $O((m+n)\log n)$
    \item $O(nm)$
    \item $O(m + n)$
\end{enumerate}

\item \textbf{(Multiple choice)} The Bellman--Ford algorithm runs in time:
\begin{enumerate}[label=(\alph*)]
    \item $O((m+n)\log n)$
    \item $O(n \log n)$
    \item $O(nm)$
    \item $O(2^n)$
\end{enumerate}

\item \textbf{(True/False)} The cut property states that the minimum-weight edge
crossing any cut belongs to some minimum spanning tree.

\item \textbf{(Multiple choice)} Kruskal's algorithm uses which data structure to
detect whether adding an edge would create a cycle?
\begin{enumerate}[label=(\alph*)]
    \item A binary min-heap
    \item A union-find (disjoint-set) structure
    \item A stack
    \item A hash map of distances
\end{enumerate}

\item \textbf{(True/False)} Prim's algorithm and Kruskal's algorithm always
produce a spanning tree of the same total weight on a connected graph.

\item \textbf{(Short answer)} In a connected graph with all distinct edge
weights, is the minimum spanning tree unique? (Yes/No, with one sentence why.)

\item \textbf{(Multiple choice)} By the cycle property, the unique heaviest edge
on a cycle belongs to:
\begin{enumerate}[label=(\alph*)]
    \item Every MST
    \item Some but not every MST
    \item No MST
    \item Exactly one MST
\end{enumerate}

\item \textbf{(Short answer)} In single-link $k$-clustering via Kruskal, you stop
when how many components remain, and what does the spacing of the clustering
equal?

\item \textbf{(True/False)} Adding the same positive constant to the weight of
\emph{every} edge of a graph can change which spanning tree is minimum.

\end{enumerate}

\newpage
\section*{Part C: Solutions}

\subsection*{Part A}

\begin{enumerate}[label=\textbf{A\arabic*.}, leftmargin=2.2em]

\item \textbf{Input:} a directed or undirected graph $G=(V,E)$, a weight $w(e)\ge
0$ on each edge, and a source vertex $s$. \textbf{Output:} for every vertex $v$,
the minimum total weight $\mathrm{dist}(v)$ of any $s$-$v$ path (or $\infty$ if
none exists). We require non-negative weights for Dijkstra because its greedy
choice -- finalize the nearest unfinished vertex and never revisit it -- is only
valid when extending a path can never \emph{decrease} its length; a negative edge
would let a longer-looking route become cheaper after the vertex is already
finalized.

\item Dijkstra maintains a set $S$ of vertices whose shortest distance is already
known, and a tentative distance $d[v]$ for each vertex (the length of the best
$s$-$v$ path found so far). It repeatedly extracts the vertex $u \notin S$ of
minimum $d[u]$, finalizes it (adds it to $S$), and \emph{relaxes} each edge
$(u,v)$: if $d[u]+w(u,v) < d[v]$, it lowers $d[v]$ to $d[u]+w(u,v)$ (recording
$u$ as $v$'s predecessor). \textbf{Invariant:} for every $u \in S$, $d[u]$ equals
the true shortest distance to $u$.

\item By induction on $|S|$. Base case: $d[s]=0=\mathrm{dist}(s)$. Inductive
step: let $u\notin S$ be the next extracted vertex, so $d[u]\le d[v]$ for all
$v\notin S$. Every relaxation sets $d[\cdot]$ to a real path length, so
$d[u]\ge\mathrm{dist}(u)$. For the reverse, take any $s$-$u$ path $P$; it leaves
$S$ at a first edge $(x,y)$ with $x\in S$, $y\notin S$. By hypothesis
$d[x]=\mathrm{dist}(x)$, and relaxing $(x,y)$ gave $d[y]\le d[x]+w(x,y)$. The
length of $P$ is at least its prefix up to $y$, which is at least
$\mathrm{dist}(x)+w(x,y)\ge d[y]\ge d[u]$ (the last step because $u$ was the
minimum). So $\mathrm{dist}(u)\ge d[u]$, hence equality. \textbf{Non-negativity
is used} at ``the length of $P$ is at least its prefix up to $y$'': discarding
the edges of $P$ after $y$ can only fail to overcount when those edges have
non-negative weight.

\item There are $n$ \textsc{ExtractMin} operations (one per vertex) and at most
$m$ \textsc{DecreaseKey}/push operations (one per edge relaxation). With a binary
heap each operation costs $O(\log n)$, so the total is $O((m+n)\log n)$, i.e.\
$O(m\log n)$ for a connected graph. (A Fibonacci heap improves this to
$O(m+n\log n)$.)

\item Vertices $s,a,b$; directed edges $s\to a$ (weight $2$), $s\to b$ (weight
$3$), $b\to a$ (weight $-2$). True $\mathrm{dist}(a)=1$ via $s\to b\to a$
($3-2$). Dijkstra extracts $a$ first with $d[a]=2$, finalizes it, and never
revisits it, so it reports $2$. The proof step that fails is the use of
non-negativity: the path $s\to b\to a$ has a prefix ($s\to b$, length $3$) that is
\emph{longer} than the whole path, because the final edge is negative.

\item Bellman--Ford relaxes all $m$ edges, $n-1$ times. The number $n-1$ is the
maximum number of edges on a simple shortest path (a path visiting $n$ vertices
has $n-1$ edges); after round $i$, $d[v]$ is correct for all vertices whose
shortest path uses at most $i$ edges, so $n-1$ rounds suffice when no negative
cycle exists. The extra ($n$-th) pass detects negative cycles: if any edge can
\emph{still} be relaxed, the graph contains a reachable negative cycle, because a
cycle-free shortest path would already have stabilized.

\item (a) Dijkstra is correct only for non-negative weights; Bellman--Ford
handles arbitrary (possibly negative) weights and detects negative cycles. (b)
Dijkstra runs in $O((m+n)\log n)$ with a binary heap; Bellman--Ford runs in
$O(nm)$. Choose Dijkstra when all weights are non-negative (it is much faster);
choose Bellman--Ford when negative edges may be present or when you must detect a
negative cycle.

\item A \textbf{spanning tree} of a connected graph $G$ is a subset of edges that
connects all $n$ vertices and contains no cycle. A \textbf{minimum spanning tree}
is a spanning tree of minimum total edge weight. A spanning tree has exactly
$n-1$ edges: a connected acyclic graph (a tree) on $n$ vertices always has
$n-1$ edges -- fewer would leave it disconnected, and one more would create a
cycle.

\item A \textbf{cut} is a partition of $V$ into two non-empty parts $(S, V
\setminus S)$. An edge \textbf{crosses} the cut if it has exactly one endpoint in
$S$ (and the other in $V\setminus S$). \textbf{Cut property:} the
minimum-weight edge crossing any cut belongs to every MST (assuming distinct
weights; with ties, it belongs to some MST).

\item Let $e=(u,v)$ be the min-weight crossing edge of cut $(S,V\setminus S)$,
$u\in S$, $v\notin S$, and suppose some MST $T$ omits $e$. Then $T\cup\{e\}$ has
exactly one cycle $C$, which contains $e$. A cycle crosses a cut an even number
of times, so $C$ has another crossing edge $e'\ne e$; since $e$ is the
\emph{minimum} crossing edge and weights are distinct, $w(e)<w(e')$. Then
$T'=(T\setminus\{e'\})\cup\{e\}$ is a spanning tree with $w(T')=w(T)-w(e')+w(e)
< w(T)$, contradicting the minimality of $T$. Hence every MST contains $e$.

\item \textbf{Cycle property:} the unique maximum-weight edge $f$ on a cycle $C$
is in no MST. \textit{Proof:} suppose MST $T$ contains $f=(x,y)$. Removing $f$
splits $T$ into two components; let $S$ be the side containing $x$, so $f$
crosses cut $(S,V\setminus S)$. Cycle $C$ also crosses this cut, so it has
another crossing edge $f'\ne f$ with $w(f')<w(f)$ (since $f$ is the unique
heaviest on $C$). Then $(T\setminus\{f\})\cup\{f'\}$ is a spanning tree of
strictly smaller weight, contradicting minimality.

\item \textbf{Kruskal:} sort edges by increasing weight; for each edge in turn,
add it iff its endpoints lie in different components (tested by union-find), and
union those components. \textit{Correctness:} when Kruskal adds edge $e=(u,v)$,
let $S$ be the set of vertices in $u$'s current component. All edges processed
before $e$ are cheaper and were either added (kept inside a component) or
discarded, so no cheaper edge crosses the cut $(S,V\setminus S)$; thus $e$ is the
minimum-weight crossing edge of this cut, and by the cut property it belongs to
every MST. So every edge Kruskal adds is safe; since it stops only when all
vertices are connected and never creates a cycle, it outputs a spanning tree of
safe edges -- the MST.

\item \textbf{Prim:} start with $S=\{r\}$ for an arbitrary $r$; repeatedly add
the minimum-weight edge with exactly one endpoint in $S$, and move that endpoint
into $S$, until $S=V$. \textit{Correctness:} at each step Prim adds the minimum
crossing edge of the cut $(S,V\setminus S)$, which is safe by the cut property;
after $n-1$ additions, $T$ is a spanning tree of safe edges. It is the
``graph analogue'' of Dijkstra because both grow a single set $S$ outward using a
heap, repeatedly pulling the cheapest frontier item -- the difference is the key:
Prim keys on the \emph{edge weight} crossing the frontier, Dijkstra on the
\emph{path distance} $d[u]=d[\text{parent}]+w$.

\item \textbf{Union-find} maintains a collection of disjoint sets. \textsc{Find}$(x)$
returns the representative of $x$'s set; \textsc{Union}$(x,y)$ merges the two
sets. \emph{Path compression} makes every node on a find-path point directly at
the root; \emph{union by rank} attaches the shorter tree under the taller. With
both, the amortized cost per operation is $O(\alpha(n))$ (inverse Ackermann --
effectively constant). In Kruskal's algorithm, each vertex starts in its own set;
an edge is added iff \textsc{Find} of its endpoints differ (different
components), after which they are \textsc{Union}ed -- this is exactly the
cycle-detection test.

\item \textbf{Uniqueness with distinct weights:} suppose $T_1\ne T_2$ are both
MSTs. Let $e$ be the minimum-weight edge in their symmetric difference (say $e\in
T_1\setminus T_2$); $e$ is unique because all weights differ. Adding $e$ to $T_2$
creates a cycle $C$; $C$ has an edge $e'\notin T_1$ (else $T_1$ would contain the
cycle), and $e'\in T_2\setminus T_1$ also lies in the symmetric difference, so
$w(e)<w(e')$ by choice of $e$. Then $(T_2\setminus\{e'\})\cup\{e\}$ is a spanning
tree lighter than $T_2$ -- contradiction. Hence the MST is unique. With equal
weights this argument breaks (the ``minimum edge in the symmetric difference''
need not be unique), and there can be several distinct MSTs of the same total
weight (e.g.\ a triangle with all equal weights has three MSTs).

\item Compare the global MST weight $W$ with two constrained quantities. (i) $e$
is in \emph{every} MST iff \emph{excluding} $e$ (computing the MST without it)
strictly increases the weight above $W$ -- equivalently, $e$ is the unique
minimum crossing edge of some cut. (ii) $e$ is in \emph{no} MST iff
\emph{including} $e$ (forcing it in) gives weight $>W$ -- equivalently, $e$ is the
unique maximum-weight edge on some cycle (cycle property). (iii) Otherwise (both
the forced-in MST and the forced-out MST achieve weight $W$) $e$ is in
\emph{some} but not every MST -- this happens with tied weights. Problem 11
implements exactly this three-way test.

\item \textbf{Single-link $k$-clustering of maximum spacing} partitions the
points into $k$ clusters so as to maximize the \emph{spacing} -- the minimum
distance between two points in different clusters. The procedure is Kruskal's
algorithm stopped once the union-find has $k$ components (equivalently, build the
MST and delete its $k-1$ heaviest edges). Stopping when $k$ components remain
maximizes the spacing; the resulting spacing equals the weight of the next edge
Kruskal would have added -- i.e.\ the $(k-1)$-th most expensive MST edge, the
smallest distance between any two of the $k$ surviving components.

\end{enumerate}

\subsection*{Part B}

\begin{enumerate}[label=\textbf{B\arabic*.}, leftmargin=2.2em]

\item \textbf{False.} Dijkstra can return incorrect distances in the presence of
negative edges, even without a negative cycle (it may finalize a vertex before
discovering a cheaper route through a negative edge). Use Bellman--Ford instead.

\item \textbf{(b) $O((m+n)\log n)$.} Each of the $n$ \textsc{ExtractMin} and up
to $m$ \textsc{DecreaseKey}/push operations costs $O(\log n)$ with a binary heap.

\item \textbf{(c) $O(nm)$.} It performs $n-1$ rounds, each relaxing all $m$
edges.

\item \textbf{True.} (And with distinct weights it belongs to \emph{every} MST;
``some MST'' is the safe statement that always holds.)

\item \textbf{(b) A union-find (disjoint-set) structure.} Adding an edge would
create a cycle exactly when its endpoints are already in the same set.

\item \textbf{True.} Both compute a minimum spanning tree, so the total weight is
the same (the minimum possible), even if they choose different edge sets when
weights are tied.

\item \textbf{Yes.} With all distinct edge weights the minimum spanning tree is
unique (see A15).

\item \textbf{(c) No MST.} This is the cycle property.

\item Stop when \textbf{$k$} components remain; the spacing equals the weight of
the next (the $(k-1)$-th heaviest MST) edge Kruskal would have added -- the
minimum distance between the $k$ resulting clusters.

\item \textbf{False.} Adding the same constant $c$ to every edge increases every
spanning tree's weight by exactly $(n-1)c$ (since each has $n-1$ edges), so the
relative order of spanning trees is unchanged and the MST stays the same.

\end{enumerate}

\end{document}
